% PREÂMBULO REORGANIZADO E COMENTADO — UTF-8 — pt-PT, AO45
% Uso: depois de documentclass, escrever input{auxiliares/preambulo.tex}.
% Motor previsto: XeLaTeX (xcolor e hyperref conservam a opção xetex).
%
% PRESERVAÇÃO E DEPENDÊNCIAS
% Todos os pacotes, opções, macros e valores do original foram conservados.
% amsmath passa a preceder unicode-math; hyperref precede o código que usa href.
% As fontes EB Garamond e NewCMMath-Book.otf têm de estar disponíveis.
% bibliografia.bib deve existir no projecto; a bibliografia requer Biber.
% logos/aeB05.pdf é usado nos cabeçalhos a partir da terceira página.
% O ambiente codigo usa number within=chapter: exige uma classe com chapter,
% como book ou report. Numa classe article, a opção tem de ser adaptada;
% conserva-se aqui a configuração original, sem mudar a numeração pretendida.
%
%=======================================================================
% 01. CORES BÁSICAS E FUNDO DE PÁGINA
%=======================================================================
% xcolor: Define cores, mistura-as e permite colorir texto, tabelas e o fundo da
% página. As opções dvipsnames e svgnames disponibilizam nomes de cores; table activa
% as cores nas tabelas.
\usepackage[xetex,dvipsnames,svgnames,table]{xcolor}

% Cores definidas por CMYK
% \definecolor{fondpaille}{cmyk}{0,0.02,0.2,0}
\definecolor{fondpaille}{cmyk}{0, 0.0118, 0.1176, 0.01} %cmyk(0%, 2.75%, 23.53%, 0%, 1%)
\definecolor{ferrugem}{cmyk}{0, 0.4457, 0.75, 0.6392}

% Cores definidas por hexadecimal
\definecolor{moleskin}{HTML}{FFF8DC}
\definecolor{laranja}{HTML}{F07241}
\definecolor{marron}{HTML}{800000}
\definecolor{lovely}{HTML}{C04848}
\definecolor{click}{HTML}{480048}
\definecolor{tinta-velha}{HTML}{4e4643}
\definecolor{Browned_Cherry}{HTML}{910831}
\definecolor{azul}{HTML}{000080}

% Cores da capa
\definecolor{BLACK}{HTML}{000000}
\definecolor{ELECTRICRED}{HTML}{DD0000}
\definecolor{TANGERINYELLOW}{HTML}{FFCE00}

% Cor consolidada para texto principal
\colorlet{cor-da-tinta}{ferrugem}
\newif\ifferrugem
\ferrugemfalse

% Cor de fundo da página
\pagecolor{fondpaille}

%=======================================================================
% 02. FONTES
%=======================================================================
% fontspec: Selecciona fontes OpenType e as suas características tipográficas com
% XeLaTeX ou LuaLaTeX.
\usepackage{fontspec} % XeLaTeX/LuaLaTeX fonts
% Mantêm-se todos os recursos OpenType originais: algarismos antigos,
% ligaduras e variantes. O sinal + activa um recurso; o sinal - desactiva-o.
\setmainfont{EBGaramond}[ % Tudo junto, não 'EB Garamond'
  %%  %%  %%  %%  %%  %%  %%  %%  %%
  Path=./auxiliares/ebgaramond/,
  Extension=.ttf,
  UprightFont=*-Regular,
  BoldFont=*-Bold,
  ItalicFont=*-Italic,
  BoldItalicFont=*-BoldItalic,
  %%  %%  %%  %%  %%  %%  %%  %%  %%
  Numbers=OldStyle,
  RawFeature={
    -aalt,
    -c2pc,
    -c2sc,
    -case,
    +dlig,
    -frac,
    -hist,
    +hlig,
    +kern,
    +liga,
    +mark,
    +mkmk,
    +onum,
    +pnum,
    +ss05,
    +cv06
  }
]
% amsmath é carregado antes de unicode-math para evitar alterações tardias.
% amsmath: Fornece alinhamentos, equações em várias linhas, intertext e outros
% recursos de composição matemática.
\usepackage{amsmath}
% unicode-math: Usa fontes matemáticas Unicode; math-style=ISO e bold-style=ISO
% escolhem as convenções ISO para os símbolos.
\usepackage[math-style=ISO, bold-style=ISO]{unicode-math}
% \setmathfont{OldStandard-Math.otf}
% \setmathfont{Garamond-Math.otf}[StylisticSet={5,6,7,9}]
\setmathfont{NewCMMath-Book.otf}

%=======================================================================
% 03. LÍNGUAS E HIFENIZAÇÃO
%=======================================================================
% polyglossia: Define as línguas, a hifenização e as designações automáticas do
% documento. A opção de língua não revê a ortografia: o texto deve ser escrito em
% pt-PT, AO45.
\usepackage{polyglossia}
\setdefaultlanguage{portuguese} %pt-PT, AO45
\PolyglossiaSetup{portuguese}{indentfirst=true}
\setotherlanguages{english,french}  % Apenas os que precisas

%=======================================================================
% 04. LAYOUT DE PÁGINA
%=======================================================================
% geometry: Define o formato e as margens da página; includehead e includefoot incluem
% cabeçalho e rodapé no cálculo da área.
\usepackage[a4paper,left=2.5cm,right=2.5cm,top=2.5cm,bottom=2.5cm,includehead,includefoot]{geometry}
% setspace: Disponibiliza comandos para alterar o espaçamento entre linhas.
\usepackage{setspace} % espaçamento de linhas
% lettrine: Cria capitulares que ocupam várias linhas no início de um parágrafo.
% calligra: Fornece uma fonte caligráfica, usada aqui nas capitulares.
\usepackage{lettrine,calligra} % letras capitulares
\renewcommand{\LettrineFontHook}{\calligra}
\LettrineTextFont{\itshape}
\setcounter{DefaultLines}{3}

% Prevenção de viúvas e órfãos
\clubpenalty=100
\linepenalty=100
\widowpenalty=100

% Outros ajustes tipográficos
% xurl: Melhora as oportunidades de quebra de linha dos endereços URL.
\usepackage{xurl}
% caption: Configura legendas; hang usa composição suspensa, small reduz o tamanho e
% bf selecciona negrito para a etiqueta.
\usepackage[hang,small,bf]{caption}
\setlength{\captionmargin}{1.5in}
\renewcommand{\fboxrule}{2pt}

%=======================================================================
% 05. TIPOGRAFIA AVANÇADA
%=======================================================================
% microtype: Melhora a composição do texto com ajustes tipográficos suportados pelo
% motor e pelas fontes.
\usepackage{microtype}      % ajustes de tipografia
% footmisc: Modifica as notas de rodapé; bottom coloca-as no fundo da área disponível
% da página.
\usepackage[bottom]{footmisc} % pés de página alinhados ao fundo
% \linespread{1.3}            % espaçamento de linha
% \linespread{1.2/1.25/1.3}

%=======================================================================
% 06. MATEMÁTICA E SÍMBOLOS
%=======================================================================
% dsfont: Fornece letras matemáticas com traço duplo através de mathds.
% mathrsfs: Fornece um alfabeto caligráfico matemático através de mathscr.
\usepackage{dsfont,mathrsfs}
% cancel: Permite riscar termos matemáticos; makeroom reserva espaço e thicklines
% engrossa os traços.
\usepackage[makeroom,thicklines]{cancel} % para riscar
% esvect: Desenha setas de vectores através de vv; conserva-se para os documentos que
% usam esse comando.
\usepackage{esvect}       % vetores

%=======================================================================
% 07. DATAS E FORMATAÇÃO NUMÉRICA
%=======================================================================
% datetime: Formata datas e horas; conservam-se as opções originais nodayofweek e level.
\usepackage[nodayofweek,level]{datetime}
% siunitx: Compõe números e unidades; output-decimal-marker selecciona a vírgula decimal.
\usepackage[output-decimal-marker={,}]{siunitx}

%=======================================================================
% 08. PACOTES DE CÁLCULO E EXERCÍCIOS
%=======================================================================
% xlop: Apresenta operações aritméticas e cálculos em disposição tradicional.
\usepackage{xlop}      % cálculo automático

%=======================================================================
% 09. GRÁFICOS E FIGURAS
%=======================================================================
% graphicx: Insere imagens e permite alterar a sua escala, rotação e recorte.
\usepackage{graphicx}
% svg: Insere imagens SVG; a conversão pode requerer Inkscape e autorização para
% executar comandos externos.
\usepackage{svg}
% tikz: Cria desenhos vectoriais, esquemas e elementos sobrepostos à página.
\usepackage{tikz}
% tkz-fct: Acrescenta recursos para desenhar funções; a execução de alguns recursos
% pode depender de ferramentas externas.
\usepackage{tkz-fct}
% tikzpagenodes: Disponibiliza nós TikZ correspondentes às áreas da página.
\usepackage{tikzpagenodes}
% circuitikz: Desenha circuitos eléctricos e os seus componentes.
\usepackage{circuitikz}
% pgfplots: Desenha gráficos de funções e dados sobre TikZ.
% pgfplotstable: Lê, transforma e apresenta dados em tabelas, em articulação com pgfplots.
\usepackage{pgfplots,pgfplotstable}
\pgfplotsset{compat=1.18}
\usetikzlibrary{
  calc,
  positioning,
  automata,
  arrows,
  shadows,
  shapes.multipart,
  matrix,
  backgrounds,
  decorations.pathreplacing,
  fit,
  plotmarks,
  tikzmark
}
\pgfdeclarelayer{background}
\pgfdeclarelayer{foreground}
\pgfsetlayers{background,main,foreground}

%=======================================================================
% 10. TABELAS AVANÇADAS
%=======================================================================
% array: Amplia a definição e a formatação das colunas das tabelas.
% longtable: Permite tabelas com continuação automática entre páginas e cabeçalhos
% repetidos.
% booktabs: Fornece regras horizontais adequadas à composição de tabelas, como top-,
% mid- e bottomrule.
% bigstrut: Acrescenta suportes invisíveis para aumentar a altura das linhas de tabelas.
% tabularx: Ajusta as colunas X à largura total indicada; o ambiente tabularx, por si
% só, não se divide entre páginas.
% seqsplit: Permite quebrar sequências longas de caracteres através de seqsplit.
\usepackage{array,longtable,booktabs,bigstrut,tabularx,seqsplit}

%=======================================================================
% 11. UTILITÁRIOS DIVERSOS
%=======================================================================
% comment: Permite excluir blocos inteiros através do ambiente comment.
\usepackage{comment}
% enumerate: Permite personalizar as etiquetas das listas enumerate. Não fornece a
% sintaxe label= e start= do pacote enumitem.
\usepackage{enumerate}
\newcounter{lastenum}
\newcommand{\storelast}{\setcounter{lastenum}{\value{enumi}}}
\newcommand{\uselast}{\setcounter{enumi}{\value{lastenum}}}
\newcommand\CC{C\nolinebreak[4]\hspace{-.05em}\raisebox{.4ex}{\relsize{-3}{\textbf{++}}}}
% rotating: Permite rodar objectos e criar figuras ou tabelas em orientação lateral.
\usepackage{rotating}
% lastpage: Cria uma referência à última página, útil para indicar o total de páginas.
\usepackage{lastpage}
% refcount: Permite obter valores numéricos a partir de referências.
\usepackage{refcount}
% standalone: Permite incorporar documentos autónomos e aproveitar os respectivos
% conteúdos.
\usepackage{standalone}
% relsize: Altera o tamanho da letra relativamente ao tamanho actual; é usado na macro CC.
\usepackage{relsize}
% ulem: Fornece sublinhados e outros realces; normalem preserva a forma habitual de emph.
\usepackage[normalem]{ulem}

%=======================================================================
% 12. SECÇÕES E SUMÁRIO
%=======================================================================
% titlesec: Configura títulos e espaçamentos e permite definir níveis adicionais de
% secção.
\usepackage{titlesec}
% tocloft: Configura a apresentação do sumário e das listas de figuras e tabelas.
\usepackage{tocloft}
% etoolbox: Fornece comandos de programação, testes e alterações de macros.
\usepackage{etoolbox}

% Novos níveis de secção
\titleclass{\subsubsubsection}{straight}[\subsubsection]
\titleclass{\subsubsubsubsection}{straight}[\subsubsubsection]
\newcounter{subsubsubsection}[subsubsection]
\newcounter{subsubsubsubsection}[subsubsubsection]
\renewcommand{\thesubsubsubsection}{\thesubsubsection.\arabic{subsubsubsection}}
\renewcommand{\thesubsubsubsubsection}{\thesubsubsubsection.\arabic{subsubsubsubsection}}

% Formato visual
\titleformat{\subsubsubsection}{\normalfont\normalsize\bfseries}{\thesubsubsubsection}{1em}{}
\titlespacing*{\subsubsubsection}{0pt}{3.25ex plus 1ex minus .2ex}{1.5ex plus .2ex}
\titleformat{\subsubsubsubsection}{\normalfont\normalsize\bfseries}{\thesubsubsubsubsection}{1em}{}
\titlespacing*{\subsubsubsubsection}{0pt}{3.25ex plus 1ex minus .2ex}{1.5ex plus .2ex}

% Sumário detalhado
\renewcommand{\cftsecfont}{\bfseries}
\renewcommand{\cftsubsecfont}{\normalfont}
\renewcommand{\cftsubsubsecfont}{\normalfont}
\makeatletter
\renewcommand{\l@section}{\@dottedtocline{1}{1.5em}{1.5em}}
\renewcommand{\l@subsection}{\@dottedtocline{2}{2.5em}{2em}}
\renewcommand{\l@subsubsection}{\@dottedtocline{3}{4em}{2.5em}}
\newcommand{\l@subsubsubsection}{\@dottedtocline{4}{6em}{3em}}
\newcommand{\l@subsubsubsubsection}{\@dottedtocline{5}{8em}{4em}}
\providecommand*{\toclevel@subsubsubsection}{4}
\providecommand*{\toclevel@subsubsubsubsection}{5}
\makeatother
\setcounter{secnumdepth}{5}
\setcounter{tocdepth}{5}

%=======================================================================
% 13. CÓDIGO-FONTE E PYTHON
%=======================================================================
% Carregar fvextra (e a sua dependência lineno) antes de csquotes.
% Esta ordem evita o aviso de fvextra e a alteração de @parboxrestore.
% fvextra: Amplia os ambientes verbatim de fancyvrb, incluindo recursos de quebra de
% linhas.
\usepackage{fvextra}
% fancyvrb: Apresenta texto e código literalmente, sem interpretar os seus caracteres
% como comandos LaTeX.
\usepackage{fancyvrb}

% pythontex: Integra código Python no documento e permite incorporar resultados. A
% execução requer o passo externo PythonTeX, além da compilação LaTeX.
% \usepackage{pythontex}

%=======================================================================
% 14. NOTAS DE RODAPÉ E CITAÇÕES
%=======================================================================
% fvextra já foi carregado na secção anterior, antes de csquotes.
% csquotes: Fornece citações e aspas adaptadas à língua, com integração em biblatex.
\usepackage{csquotes}
% appendix: Organiza os apêndices; toc e page permitem a sua entrada no sumário e uma
% página própria.
\usepackage[toc,page]{appendix}
\renewcommand\appendixpagename{Apêndices}
% biblatex: Gere as referências e citações bibliográficas. Este documento usa Biber,
% citações numéricas compactas e a ordem de citação.
\usepackage[
  backend=biber,
  style=numeric-comp,
  language=auto,
  natbib=false,
  block=space,
  isbn=false,
  url=true,
  doi=true,
  mcite=true,
  eprint=true,
  sorting=none
]{biblatex}
\addbibresource{bibliografia.bib}

% As notas normais usam letras; a marca é apresentada entre parênteses.
% symbolfootnote permite escolher uma marca simbólica para uma nota.
% Formato das notas de rodapé
\renewcommand{\thefootnote}{\alph{footnote}}
\makeatletter
\renewcommand\@makefnmark{\@textsuperscript{\normalfont(\@thefnmark)}}
\makeatother
\long\def\symbolfootnote[#1]#2{%
  \begingroup
  \def\thefootnote{\fnsymbol{footnote}}
  \footnote[#1]{#2}
  \endgroup
}

%=======================================================================
% 15. CÓDIGO E CAIXAS (QUEBRÁVEIS ENTRE PÁGINAS)
%=======================================================================
% listings: Apresenta código-fonte com realce de sintaxe, numeração de linhas e
% estilos; não executa o código.
\usepackage{listings}
% tcolorbox: Cria caixas com títulos, cores e molduras; most carrega um conjunto amplo
% de bibliotecas.
\usepackage[most]{tcolorbox} % carregar UMA SÓ VEZ no preâmbulo
\tcbuselibrary{listings,breakable}

%=======================================================================
% 16. CORES (CÓDIGO)
%=======================================================================
\definecolor{codebg}{gray}{0.96}
\definecolor{codeframe}{gray}{0.75}
\definecolor{codeblue}{RGB}{40,40,180}
\definecolor{codegreen}{RGB}{0,120,0}
\definecolor{codered}{RGB}{160,20,20}
\definecolor{darkgray}{gray}{0.4}

%=======================================================================
% 17. TRADUÇÃO LISTING → LISTAGEM
%=======================================================================
\renewcommand{\lstlistingname}{Listagem}
\renewcommand{\lstlistlistingname}{Lista de Listagens}

%=======================================================================
% 18. ESTILO LISTINGS (ESTÁVEL COM XELATEX)
%=======================================================================
% Definir o estilo não o aplica a todos os ambientes lstlisting.
% Usar [style=codigo] num lstlisting ou o ambiente codigo definido abaixo.
\lstdefinestyle{codigo}{
  language=Python,
  basicstyle=\ttfamily\small,
  keywordstyle=\color{codeblue}\bfseries,
  commentstyle=\color{codegreen}\itshape,
  stringstyle=\color{codered},
  numbers=left,
  numberstyle=\tiny\color{darkgray},
  numbersep=8pt,
  breaklines=true,
  showstringspaces=false,
  columns=fullflexible,
  keepspaces=true
}

% ---------------------------------------------------------
% Ambiente de código quebrável
% Uso: \begin{codigo}{Título} ... \end{codigo}
% ---------------------------------------------------------
\newtcblisting[auto counter, number within=chapter]{codigo}[2][]{
  breakable,
  enhanced,
  listing only,
  listing options={style=codigo},
  colback=codebg,
  colframe=codeframe,
  boxrule=0.4pt,
  arc=2pt,
  left=6mm,right=2mm,top=2mm,bottom=2mm,
  title={Listagem~\thetcbcounter: #2},
  fonttitle=\bfseries,
  #1
}

%=======================================================================
% 19. CAIXA PARA PROMPTS CIENTÍFICAS / ORIENTAÇÕES IA (VERBATIM)
%=======================================================================
\newtcblisting{promptcientifico}[2][]{%
  breakable,
  enhanced,
  listing only,
  colback=gray!10,
  colframe=black,
  boxrule=0.6pt,
  arc=2pt,
  left=4mm,
  right=4mm,
  top=2mm,
  bottom=2mm,
  title={#2},
  fonttitle=\bfseries,
  listing options={
    basicstyle=\ttfamily\small,
    breaklines=true,
    columns=fullflexible,
    keepspaces=true,
    showstringspaces=false
  },
  #1
}



%=======================================================================
% 20. HYPERREF
%=======================================================================
% hyperref: Cria ligações e marcadores no PDF. Conservam-se as opções originais; o
% último colorlinks=true prevalece sobre colorlinks=false.
\usepackage[%
  bookmarks=true,
  citecolor=black,
  filecolor=black,
  linkcolor=black,
  urlcolor=blue,     % Cor específica para URLs
  unicode=false,
  colorlinks=false,
  plainpages=false,
  pdfpagelabels,
  colorlinks=true,  % Ativa cores em vez de caixas
  xetex]{hyperref}

%=======================================================================
% 21. TOPO E RODAPÉ DE PÁGINAS (LOGOS E TEXTO)
%=======================================================================
% eso-pic: Adiciona elementos ao fundo ou ao primeiro plano das páginas no momento de
% as enviar para o PDF.
\usepackage{eso-pic}

\AddToShipoutPictureBG{%           % Início do bloco que adiciona elementos ao fundo da página
  \ifnum\value{page}>2           % Condição: só aplica a partir da página 3
    \AtPageUpperLeft{%            % Posiciona tudo relativamente ao canto superior esquerdo da página
      \begin{tikzpicture}[remember picture, overlay]% Início do ambiente TikZ sobreposto (overlay) e lembrando posição
        % Linhas superior e inferior
        \draw[line width=0.5pt, color=black] (2.5cm, -3cm) -- (\paperwidth - 2.5cm, -3cm); % Linha horizontal superior
        \draw[line width=0.5pt, color=black] (2.5cm, -\paperheight + 3cm) -- (\paperwidth - 2.5cm, -\paperheight + 3cm); % Linha horizontal inferior
        
        % Conteúdo página ímpar
        \ifodd\value{page}           % Verifica se a página é ímpar
          \node [anchor=west, text depth=0pt] at (2.5cm, -2.7cm){\tiny\href{mailto:luisav.ferreira@gmail.com}{Grupo Disciplinar 510}}; % Texto canto superior esquerdo
          \node [anchor=north east] at (\paperwidth-2.5cm, -2.2cm){\includegraphics[width=0.05\paperwidth]{logos/aeB05.pdf}}; % Logo canto superior direito
          \node [anchor=west, text height=0pt] at (2.5cm, -\paperheight + 2.7cm){\tiny\href{mailto:luisav.ferreira@gmail.com}{Luís A.V. Ferreira}}; % Texto canto inferior esquerdo
          \node [anchor=west] at (2.5cm, -\paperheight + 2.5cm){\tiny\jobname.tex \smash{\rule[-0.5ex]{0.6pt}{1ex}}}; % Nome do ficheiro e linha decorativa
        % Conteúdo página par
        \else                        % Caso seja página par
          \node [anchor=west, text depth=0pt] at (2.5cm, -2.7cm){\includegraphics[width=0.05\paperwidth]{logos/aeB05.pdf}}; % Logo canto superior esquerdo
          \node [anchor=east, text depth=0pt] at (\paperwidth - 2.5cm, -2.7cm){\tiny\href{mailto:luisav.ferreira@gmail.com}{Grupo Disciplinar 510}}; % Texto canto superior direito
          \node [anchor=east, text height=0pt] at (\paperwidth - 2.5cm, -\paperheight + 2.7cm){\tiny\href{mailto:luisav.ferreira@gmail.com}{Luís A.V. Ferreira}}; % Texto canto inferior direito
          \node [anchor=east] at (\paperwidth - 2.5cm, -\paperheight + 2.5cm){\tiny\jobname.tex \smash{\rule[-0.5ex]{0.6pt}{1ex}}}; % Nome do ficheiro e linha decorativa
        \fi
      \end{tikzpicture}
    }
  \fi
}

